% !PRINTED 277
\weberpage{277}
In agreement with this, $\sqrt{-2}$ is likewise real or purely imaginary
(Bd.\,I, \S.\,138, 4., 6.), and consequently in these cases $\Theta$,
$\psi$, $\chi$, and hence the whole tetrahedral group, are real.

2) But if $p \equiv 3 \pmod 4$, then $\gamma^{\frac{p^2-1}{4}}$,
$\gamma^{-\frac{p^2-1}{4}}$ are conjugate imaginary, because
$\gamma^{(p+1)\frac{p^2-1}{4}} = 1$.

If $p \equiv 3 \pmod 8$, then $\sqrt{-2}$ is real and
$\gamma^{\frac{p^2-1}{8}(p+1)} = -1$, hence $\gamma^{\frac{p^2-1}{8}}$
and $-\gamma^{\frac{p^2-1}{8}}$ are conjugate imaginary; and if $p \equiv 7
\pmod 8$, then $\sqrt{-2}$ is purely imaginary, $\gamma^{\frac{p^2-1}{8}}$
and $\gamma^{-\frac{p^2-1}{8}}$ are conjugate imaginary, and consequently
$\Theta$, $\psi_1$, $\chi$ belong in these cases to the group $\Gamma_p$.
There is thus in \gesperrt{all cases in the group} $L_p$ a tetrahedral group.

For $p = 5$, for example, since $-2$ is a quadratic non-residue of 5,
one may take $\varepsilon = \sqrt{-2}$, and, since $\pm 2$ are the two
primitive roots of 5, say
\[
\gamma^6 = -2, \quad \gamma^3 = \sqrt{-2}, \quad \gamma = 1 + \sqrt{-2}
\]
then it follows that
\weq{6}{\begin{aligned}
&\Theta = \begin{pmatrix} 2, & 0 \\ 0, & 3 \end{pmatrix},\quad
\psi = \begin{pmatrix} 0, & 2 \\ 2, & 0 \end{pmatrix},\quad
\psi_1 = \begin{pmatrix} 0, & 1 \\ -1, & 0 \end{pmatrix} \\
&\chi = \begin{pmatrix} 1, & 1 \\ 2, & -2 \end{pmatrix},\quad
\chi^2 = \begin{pmatrix} 2, & 1 \\ 2, & -1 \end{pmatrix},
\end{aligned}}
and from this results the tetrahedral group sought for $p = 5$
[\S.\,56, (3)]:
\weq{7}{\begin{aligned}
&\begin{pmatrix} 1, & 0 \\ 0, & 1 \end{pmatrix},
\begin{pmatrix} 2, & 0 \\ 0, & -2 \end{pmatrix},
\begin{pmatrix} 0, & 2 \\ 2, & 0 \end{pmatrix},
\begin{pmatrix} 0, & 1 \\ -1, & 0 \end{pmatrix} \\
&\begin{pmatrix} 1, & 1 \\ 2, & -2 \end{pmatrix},
\begin{pmatrix} 2, & 2 \\ 1, & -1 \end{pmatrix},
\begin{pmatrix} -1, & 1 \\ 2, & 2 \end{pmatrix},
\begin{pmatrix} -2, & 2 \\ 1, & 1 \end{pmatrix} \\
&\begin{pmatrix} 2, & 1 \\ 2, & -1 \end{pmatrix},
\begin{pmatrix} -1, & 2 \\ 1, & 2 \end{pmatrix},
\begin{pmatrix} 1, & 2 \\ 1, & -2 \end{pmatrix},
\begin{pmatrix} -2, & 1 \\ 2, & 1 \end{pmatrix}.
\end{aligned}}
This group is a divisor of $L_5$ of index 5.

IV. \gesperrt{Octahedral group.}

In the octahedral group a cyclic group of degree 4 is contained, and
such a group can therefore exist only when $1/2\,(p-1)$ or
$1/2\,(p+1)$ is divisible by 4, that is, when $p \equiv \pm 1 \pmod 8$.
